\documentclass[../Main.tex]{subfiles}

\begin{document}

\section{Evaluation} \label{sec:Evaluation}

\subsection{Ziel des Playtests}

\begin{itemize}[leftmargin=1.5cm]
    \item \textcolor{red}{Aus Recherche (aus Abschnitt~\ref{sec:Prototyping} verschoben): Der geplante Playtest entspricht dem späten Typ (Fehler und Politur); ein früher Low-Fidelity-Test der Gegnerplatzierung fehlt bislang.}
\end{itemize}

\subsection{Aufbau des Playtests}

\begin{itemize}[leftmargin=1.5cm]
    \item \textcolor{red}{Aus Recherche (aus Abschnitt~\ref{sec:Prototyping} verschoben): breites, informelles Testen (auch mit unerfahrenen Personen) deckt Bedienprobleme auf, die Entwickler übersehen (HangLine); Feedback gefiltert, aber ernst nehmen (Taylor, Kitchen-Metapher).}
\end{itemize}

\subsubsection{Testpersonen}

\subsubsection{Fragebogen}

\subsubsection{Beobachtungsraster}

\subsection{Durchführung}

\begin{itemize}[leftmargin=1.5cm]
    \item \textcolor{red}{Aus Recherche (aus Abschnitt~\ref{sec:Prototyping} verschoben): Playtest-Regeln (Totten): nicht eingreifen oder erklären, Happy Accidents zulassen, Feedback phasengerecht einholen, Redundanzregel für Signale (mehrere verschiedenartige Hinweise gleichzeitig).}
\end{itemize}

\subsection{Ergebnisse}

\subsubsection{Quantitative Ergebnisse}

\subsubsection{Qualitative Beobachtung}

\subsection{Diskussion}

\subsubsection{Interpretation der Ergebnisse}

\subsubsection{Verbesserungsvorschläge}

\end{document}